\documentclass{turabian-researchpaper}
\usepackage[utf8]{inputenc}
\usepackage[english]{babel}
\usepackage[utf8]{inputenc}
\usepackage{csquotes, ellipsis}
\usepackage[authordate-trad,backend=biber]{biblatex-chicago}
\usepackage{outlines}
\bibliography{quantrdbib.bib}
\title{Put Your Title Here}
\author{Put your name here}
\date{\today}
\raggedright

\begin{document}
\maketitle
\section{General Problem Area}
This is where you hook the reader. Show why your paper is interesting. make sure readers will want to know the answer to your research question because you have convinced them it is important. You can start to narrow. 

\section{Research Question}
Tell the reader your specific research question. Make sure your question is clear, testable, narrow, relevant to political policy, significant to political theory, and original. Your research question should clearly have a dependent variable and an independent variable. make sure your question is clear, testable, sufficiently narrow, practically relevant, theoretically significant, and original. Also make sure that your question is general enough that you will have more than one case. This section only needs to be a sentence or two.

\section{Literature Review}
Here is where you sum up previous research that has been done that relates to your topic. Don't just write one paragraph about each of five different sources. Organize the previous research. Help the reader to see the big picture. Where does your research fit in to what others have already done? How have others explained your dependent variable? What has been left out of previous research? 

Do not forget that one of the best reasons for using \LaTeX{} is that it will create your bibliography for you. If you don't know how to do citation and bibliography, so back to the document called ``\LaTeX{} for POSC 300 Students.'' All you need to do is enter the proper citations in this document, and copy and paste the bibliography entery from Google Scholar into the document called ``quantrdbib.bib'' and \LaTeX{} will make the perfect Turabian reference page for you.

\section{Theory}
What is your causal story? What is the logic that gets from your independent variable to your dependent variable? What are the causal mechanisms? Note that while the literature review was a summary of other people's arguments, the theory section is YOUR argument. Take your time and do a good job here. This section will take you more than a page. It may take several pages.

\section{Hypothesis}
This section needs to lay out at least one hypothesis. Remember that a hypothesis is an observable implication of your theory. Make sure there are clear independent and dependent variables, and that those variables are specific enough that you can make a good operationalization for them (which will happen in the next section).

When I write papers, I like to set my hypotheses off, and you can do the same thing. Here is the format I use to make my hypotheses stand out:

I believe that I have the supernatural ability to affect the outcome of football games by cheering for a particular team. As I explained in the theory section of this paper, when I cheer for a football team, they are more likely to lose the game. Therefore, my first hypothesis is:

\begin{description}
		\item \textbf{Hypothesis}: there will be a positive and significant relationship between my cheering for a specific football team and that team losing the game, \emph{ceteris paribus}.
	\end{description} 

You can use this same format and just put your own hypothesis into code I have written above.

\section{Operational Definitions}
In this section, you need to give an operational definition for your dependent and independent variables. These may be your own definitions, or an explanation of why you are using another researcher's definition (and the appropriate citation). 

\section{Control Variables}
In this section, you need to discuss that variables you will control for in your research. These need to be other important causes of your dependent variable (other than your independent variable). Use your literature review to identify these important variables. You should have at least three or four control variables for this paper.

\section{Case Selection}
This section of the paper needs to be very different from your Qualitative RD. When using a quantitative research design, we use many cases (as many as possible/appropriate). Thinking about your topic, decide what the population looks like. If you could use the whole population, they do so. For example, if your unit of analysis is countries, you can just use all countries. If you cannot use the whole population, then you need to use a sample.For example, if your unit of analysis is individuals (living in Australia), you will have to use a sample (instead of getting information from each person living in Australia). How will you choose which cases belong in your sample?

\section{Data and Criteria for Verification}
Identify what kinds of information you will need to find to make your case studies. Be specific. Don't say ``I will need to find books and articles on the effects of business interests on coups in Fiji and Argentina.'' Say something more like, ``I will need to find evidence of direct meetings of business leaders with people who would carry out coups in Fiji and Argentina. I will also need to find evidence that the people who participated in the coup received some kind of financial reward from business interests after the coup occurred.'' You should be looking for places to go to get datasets for this paper.

After you describe the kinds of information you would need to find, discuss what kinds of evidence would give support to your hypothesis, and what kinds of information would refute your hypothesis. Remember to talk about what kinds of statistical results you will look for in this paper.

\section{Data Sources}
Here you should identify at least five specific data sources that you could go to in order to find the kind of information you described in the previous section. These sources will be different from the sources you used in your literature review. These will be sources that you use to find the specific kind of evidence you are looking for. Offer a brief discussion of how each of these sources will be useful to you. Do not forget to include each of these sources on your reference page.

\section{Limitations}
In this section, you will describe the weaknesses of your project. All research has weaknesses. Just be honest about yours.

\section{Tenative Outline}
Here is the \LaTeX{} format to create an outline. All you need to do is put your own points into this outline. You should be able to figure out how to create more or less subpoints fairly easily, but ask if you need help.

\begin{outline}
\1 General problem area
    \2 a point about your GPA
    \2 another point about your GPA
\1 Research question
    \2 a point about your research question
        \3 a sub point about your research question
\1 Just continue on in this same pattern
\end{outline}






\clearpage
\printbibliography
\end{document}
